% --- Patron: Morag the Hag, Weaver of Hidden Costs (Faerie Bargains & Hidden Costs) ---
\subsection{Morag the Hag, Weaver of Hidden Costs (Faerie Bargains \& Hidden Costs)}

\textit{Lore.} In the hour between twilight and dawn, where borders thin and mortal cunning meets older cleverness, Morag sets her webs of bargain and consequence. She is the grandmother who offers exactly what you need, the forest crone with a perfect solution, the crossroads dweller who makes dreams come true—for a price that only reveals itself once the road bends back. Her wisdom is simple and sharp: every gift carries obligation, every kindness creates debt, and every wish must be paid in coin not named aloud.

Her bargains are never unfair—only incomplete. She grants what is promised; the true cost lies in what is unspoken, written in thorn‑scratch, and due when the moon has turned thrice.

\begin{quote}
``She gives you the gold, but not the weight it carries. She grants your heart's desire, but not the hunger it awakens. She answers your plea, but not the price that pleases her. In Morag's court, every `yes' is a `yes, but\ldots\,' ''
\end{quote}

\noindent\textbf{Domain Focus}
\begin{itemize}
  \item \textbf{Hidden Obligations:} reading the unspoken costs in every bargain
  \item \textbf{Faerie Cunning:} strict courtesy, sharp clauses, old‑road law
  \item \textbf{Hearth‑Way Wisdom:} folk practice for spirits, thresholds, and natural justice
  \item \textbf{Consequential Magic:} power that always arrives with strings attached
\end{itemize}

\subsubsection*{Rites of Morag}

\paragraph{Rite of the Unfinished Promise (Low, 4 XP)} \hfill \\
\emph{Scene; Near; No.} \\
\textbf{Materials:} A half‑written contract or a promise spoken but not sealed. \\
\textbf{Effect:} Frame a bargain with an escape clause only you understand. Gain \textbf{+1 die} when the other party acts to fulfil their obligation; you may interpret ambiguous terms in your favour. \\
\textbf{Invoke:} 1 action; mark \textbf{+1 Obligation}. \\
\emph{Requires: Familiar (\textit{Invoke:} 1 Boon).}

\paragraph{Rite of the Borrowed Luck (Low, 5 XP)} \hfill \\
\emph{Scene; Touch; Yes.} \\
\textbf{Materials:} A token of good fortune from someone who recently succeeded. \\
\textbf{Effect:} Share in another’s fresh success. Gain \textbf{+2 dice} on one roll that benefits from their good fortune. \\
\textbf{Push It:} The luck spreads but creates debt – the original bearer suffers misfortune later; mark \textbf{1 SB (Hearts)}. \\
\textbf{Invoke:} 1 action; mark \textbf{+1 Obligation}. \\
\emph{Requires: Familiar (\textit{Invoke:} 1 Boon).}

\paragraph{Rite of the Speaking Knot (Standard, 8 XP)} \hfill \\
\emph{Scene; Near; Yes.} \\
\textbf{Materials:} Cord or thread tied with counted knots while speaking the bargain. \\
\textbf{Effect:} Bind a promise with fae knots that tighten when broken. Both parties gain \textbf{+1 die} while keeping the bargain; breaking it escalates consequences (Harm 1, then Harm 2, then \textit{Curse}). \\
\textbf{Push It:} The knots grow willful and seek payment from the breaker’s kin; mark \textbf{+1 Obligation}. \\
\textbf{Invoke:} 1 action; mark \textbf{+1 Obligation}. \\
\emph{Requires: Familiar + Codex (\textit{Invoke:} 1 Boon).}

\paragraph{Rite of the Threefold Exchange (Standard, 7 XP)} \hfill \\
\emph{Scene; Near; Yes.} \\
\textbf{Materials:} Three items of equal seeming value but different true worth. \\
\textbf{Effect:} Broker a trade where each party yields hidden value. All participants gain \textbf{+1 die} when the exchange benefits them; Morag claims the difference in true worth as her tithe. \\
\textbf{Push It:} The exchange binds across generations; mark \textbf{1 SB (Spades)} as ancestral debts stir. \\
\textbf{Invoke:} 1 action; mark \textbf{+1 Obligation}. \\
\emph{Requires: Familiar + Codex (\textit{Invoke:} 1 Boon).}

\paragraph{Rite of the Cauldron’s Secret (High, 13 XP)} \hfill \\
\emph{Extended; Zone; Yes.} \\
\textbf{Materials:} A cauldron never yet empty, ingredients dearer than coin, and a willing gift of time. \\
\textbf{Effect:} Brew a solution to any problem, but the cauldron asks one more ingredient – one the brewer does not realise they are providing until too late. \\
\textbf{Push It:} The remedy works perfectly but breeds dependence – users must return or suffer withdrawal; mark \textbf{+2 Obligation}. \\
\textbf{Invoke:} Extended rite; mark \textbf{+2 Obligation}. \\
\emph{Requires: Familiar + Codex + Tier III (\textit{Invoke:} 2 Boons).} \\
\emph{Obligation:} 7 segments.

\paragraph{Rite of the Crossroads Court (High, 14 XP)} \hfill \\
\emph{Extended; Zone; Yes.} \\
\textbf{Materials:} A true crossroads at twilight, offerings set to the four ways, and the tears of one who bargained thrice. \\
\textbf{Effect:} Convene Morag’s court: all present must bargain. Each pact grants a marked boon and hides a cost that will ripen at the worst time. \\
\textbf{Push It:} Judgments bind across realms and seasons; mark \textbf{+2 Obligation} and begin \textbf{Entangled Fates [6]}. \\
\textbf{Invoke:} Extended rite; mark \textbf{+2 Obligation}. \\
\emph{Requires: Familiar + Codex + Tier III (\textit{Invoke:} 2 Boons).} \\
\emph{Obligation:} 8 segments.

\subsubsection*{Cantors \& Cults of Morag}

\textbf{The Cantor’s Song.} A Cantor of Morag sings in \textit{harmonies of the half‑heard} – melodies that seem to come from just over your shoulder, from a room you left, from a memory you forgot. Their voice is used to sweeten a bitter bargain, to make a debtor forget the interest rate until it is too late, or to call a crossroads spirit to witness an oath. Morag’s Cantors are often found in marketplaces, singing the fine print into existence.

\textbf{The Cult of the Tarnished Cauldron.} Morag’s cults meet at crossroads, in root cellars, and around cauldrons that never cool. The Tarnished Cauldron is a fellowship of hedges‑witches, debt‑collectors, and those who have learned that every gift is a hook. A ritual of the Cult involves each member adding an ingredient to the cauldron (a hair, a coin, a written wish) while speaking a price they are willing to pay – a price that the cauldron remembers. Their creed is written in chalk on the cauldron’s belly: \textit{“What is freely given can never be reclaimed. What is bargained can always be renegotiated.”} Outsiders are welcome to bring a gift and ask a boon – but they will leave with a debt.

\subsubsection*{Witchcraft of Morag (Craft / Hag’s Work)}

A witch who walks with Morag is a \textit{knot‑wife} or \textit{cauldron‑tender} – one who crafts bargains into being with thread, iron, and simmering pot. Their magic works through the tools of the hedge: a needle that sews a promise into a hem, a pair of shears that cut an obligation cleanly (or not), or a ladle that stirs consequence into soup. In Aeler, a True‑Mason might carve a “bargain stone” at a bridge’s midpoint, so that every crossing is a silent contract. In Aelaerem, a hedgewitch might tie a red thread around a debtor’s gatepost, binding them until the debt is paid.

\textbf{The Witch’s Tool: The Bargain Needle.} A needle of cold iron, its eye shaped like a knot. Once per session, the witch may sew a “bargain seam” into a garment, a threshold, or a document. The seam is invisible to all but those who owe a debt to Morag. When the bargain is fulfilled (or broken), the thread loosens (or tightens). The witch may also use the needle to “unpick” a broken promise, resetting a failed negotiation clock by 1 segment at the cost of 1 Fatigue.

\textbf{A Hedge Gift: Cauldron’s Whisper (2 XP).} Once per scene, the witch may stir any pot or bowl of liquid (water, broth, wine) with a finger and speak a name. For the next hour, the liquid will “remember” any promise or debt spoken over it. Drinking from the liquid grants the drinker \textbf{+1 die} to recall the terms of a bargain, or to detect if someone is lying about an obligation. The liquid tastes faintly of iron and hawthorn.

\subsubsection*{Morag’s Corruption}

\begin{tblr}{
  colspec = {Q[c,1cm] X[2.5] X[2.5]},
  rowhead = 1,
  row{odd} = {bg=gray!10},
  row{1} = {bg=gray!30, font=\bfseries},
  hlines,
}
Tier & Benefit & Cost / Quirk \\
1 & \textbf{Faerie Sight:} +1 die to \textit{Notice} hidden obligations or unspoken costs. & \textbf{Bargain Sensitivity:} –1 die when refusing offers or walking away; the urge to negotiate bites. \\
2 & \textbf{Cunning Exchange:} Once/scene, propose a trade where you gain equal or greater value for lesser seeming payment. & \textbf{Debt Awareness:} Ever‑counting ledgers of owed/owing; suffer \textbf{Fatigue 1} in crowded markets. \\
3 & \textbf{Old‑Road Wisdom:} +2 dice when dealing with spirits, fae, or ancient powers. & \textbf{Obligation Magnet:} “Fair trades” find you – each with costs not yet clear. \\
4 & \textbf{Bargain Mastery:} Once/session, set terms so favourable that even the other party feels satisfied (for now). & \textbf{Price Collector:} Collect on at least one outstanding debt each session or suffer –1 die to social rolls. \\
5 & \textbf{Crossroads Power:} Once/session, call on a crossroads where you bargained before. & \textbf{Fate Entanglement:} Your pacts snag your circles; –1 die to keep agreements purely personal. \\
6+ & \textbf{Hag’s Court:} Once/session, arbitrate any bargain; your terms bind all parties. & \textbf{Debt Incarnate:} Mark +3 Obligation; risk becoming a walking covenant, unable to act without minting new debts. \\
\end{tblr}

\subsection*{Playstyle Notes}
Morag favours negotiators who read the \emph{said} and the \emph{meant} at once. She rewards courtesy kept, clauses counted, and solutions bought with clever coin. The cost is a life lived on ledgers: every boon with a braid, every kindness with a key.

\noindent\textbf{Emphasises}
\begin{itemize}
  \item \textbf{Faerie Bargaining:} strict forms, sharp prices, safe exits
  \item \textbf{Hidden Costs:} the string on every gift
  \item \textbf{Hearth‑Way Wisdom:} spirits, thresholds, and natural justice
  \item \textbf{Consequential Magic:} boons that demand upkeep
  \item \textbf{Bargain Mastery:} mutually pleasing terms that still pay you twice
\end{itemize}